\documentclass[10pt,a4paper]{amsart}
\usepackage[margin=19mm]{geometry}
\usepackage{amsmath,amssymb,mathtools}
\usepackage[scr=rsfs,cal=euler]{mathalfa}
\usepackage{xfrac}
\usepackage[unicode,hidelinks,bookmarks=false]{hyperref}
\newcommand{\ie}{i.e.\ }
\newcommand{\cf}{cf.\ }
\newcommand{\resp}{respectively\ }
\newcommand{\Subsection}{\subsection}
\providecommand{\nobreakdash}{\nobreak\hbox{-}}
\setlength{\emergencystretch}{2em}
\multlinegap0pt
\usepackage{tikz-cd}
\usepackage{xparse}
\usepackage{enumitem}
\newcommand{\category}[1]{\mathcal{#1}}
\newcommand{\A}{\category{A}}
\newcommand{\B}{\category{B}}
\newcommand{\C}{\category{C}}
\newcommand{\D}{\category{D}}
\newcommand{\E}{\category{E}}
\newcommand{\G}{\category{G}}
\NewDocumentCommand{\Fun}{O{}D<>{}mm}{\operatorname{Fun}_{#1}^{\mathrm{#2}}(#3,#4)}
\theoremstyle{plain}
\newtheorem{theorem}{Theorem}[subsection]
\newtheorem{proposition}{Proposition}[subsection]
\newtheorem{lemma}{Lemma}[subsection]
\newtheorem{corollary}{Corollary}[subsection]
\theoremstyle{definition}
\newtheorem{definition}{Definition}[subsection]
\theoremstyle{remark}
\newtheorem{remark}{Remark}[subsection]
\title{Stable $\infty$-categories for representation theorists}
\author{Gustavo Jasso}
\date{Source: arXiv:2603.14970v3 (CC BY 4.0); the author's own native \LaTeX{} source, set here in the editorial AMS \texttt{amsart} wrapper (the source's own class is \texttt{amsart} a4paper).}
\begin{document}
\maketitle
\noindent\textit{Editorial scope.} The counted claim of this unit is the article's Proposition 2.9.7 (source label \texttt{prop:lax-limit}), the self-referential universal property of the upper-triangular gluing as a lax limit, delivered together with the author's complete original proof. The statement and the proof are the only retained passages and every one of their characters is the author's own; the proof is delivered as the single primary proof body exactly as the author's proof environment encloses it. Printed numbering: the article declares its theorem-like environments inside its own macro file, which is not redistributed, so the wrapper restates the same declarations, and the section, subsection and proposition counters are advanced so that the delivered PDF prints the counted claim as Proposition~2.9.7, exactly as the published article does. Omitted: the rest of the article, that is its title block, abstract, introduction, the primer on infinity-category theory, Section~2 up to the counted proposition, and the sections that follow it. The retained unit is the maximal source-local passage of the article that is closed under its own cross-references: the surrounding text of this proposition refers to definitions and to results of the article whose own statements and proofs use further material that cannot be retained inside a unit of this size, so nothing else is included. Nothing of the counted statement or of the counted proof is omitted. Class substitution: the source class is the AMS \texttt{amsart} class with the author's own packages and his own macro file; the wrapper is AMS \texttt{amsart} carrying the author's own macro declarations, including his document-command for the exact functor category and his calligraphic category letters. Reference handling: no retained passage contains a citation; the article's bibliography is a separate file that is kept verbatim in the eprint directory and is not redistributed, and no reference entry is transcribed, delivered or invented. No mathematics is written, repaired or re-derived here.

\setcounter{section}{2}
\setcounter{subsection}{9}
\setcounter{proposition}{6}
\begin{proposition}
  \label{prop:lax-limit}
  Let $F\colon\B\to\A$ be an exact functor between stable $\infty$-categories.  For each stable $\infty$-category $\C$, there is a canonical equivalence of
  stable $\infty$-categories
  \[
    \Fun<ex>{\C}{\B\vec{\times}_F\A}\stackrel{\sim}{\longrightarrow}\Fun<ex>{\C}{\B}\vec{\times}_{F_*}\Fun<ex>{\C}{\A},\quad\Phi\longmapsto(p_*(\Phi),(i_R)_*(\varepsilon_\Phi)),
  \]
  where $F_*$ denotes the post-composition functor
  \[
    F_*\colon\Fun<ex>{\C}{\B}\longrightarrow\Fun<ex>{\C}{\A},\qquad
    \Phi\longmapsto F\circ\Phi,
  \]
  and where $\varepsilon\colon p_Lp\to\mathbf{1}_{\C}$ is the counit of the adjunction
  $p_L\dashv p$.
\end{proposition}

\begin{proof}
  The existence of the claimed equivalence follows immediately from the
  existence of the induced cartesian square
  \[
    \begin{tikzcd}
      \Fun<ex>{\C}{\B\vec{\times}_F\A}\rar\dar[swap]{p_*}\ar[phantom]{dr}[description]{\text{PB}}&\Fun{s\to t}{\Fun<ex>{\C}{\A}}\dar{s^*}\\
      \Fun<ex>{\C}{\B}\rar[swap]{F_*}&\Fun<ex>{\C}{\A},
    \end{tikzcd}
  \]
  where we implicitly use the canonical identification
  \[
    \Fun<ex>{\C}{\Fun{s\to t}{\A}}\simeq\Fun{s\to t}{\Fun<ex>{\C}{\A}}.\qedhere
  \]
\end{proof}

\end{document}
